\subsection{\problemblock{*Title} Data Block}\label{sec:block-title}

A problem title consisting of one or more lines of text is optional. When supplied,
it is printed at the top of every page of standard output and at the beginning of EGS
database files. A title helps identify and document a problem. There is no formal
limit on the number of title lines, although printout becomes cluttered with more
than five or six.

\begin{lstlisting}[style=taosmanual]
*title This is the first line of the title
       this is the second line of the title
       and so on, until a data block name is hit

*atmos standard
\end{lstlisting}

The title begins with the first nonblank character after the
\problemblock{*title} keyword and extends to the next valid data-block keyword. In
the example, it begins at \texttt{This} and ends immediately before
\problemblock{*atmos}. A title can therefore contain arbitrary text except a valid
data-block keyword. Because all such keywords begin with an asterisk, difficulties
can be avoided by not placing an asterisk at the beginning of a word in title text.
